\section{O Axioma do Par}\label{sec:axiomaDoPar}
\index{axioma!do par}
\index{par!não ordenado}
\index{par!ordenado}

Por enquanto, os únicos axiomas introduzidos foram o Axioma da Extensão e o Esquema da Compreensão.
O primeiro é fundamental para a caracterização da igualdade entre conjuntos, enquanto o segundo é crucial para a construção, a partir de um conjunto já existente, de subconjuntos definidos por fórmulas.

Porém, juntos, esses axiomas não garantem a existência de nenhum conjunto particular além do vazio.

Para introduzir novos conjuntos, precisamos de novos axiomas.

O primeiro deles é o Axioma do Par.

Em sua axiomatização de 1908, geralmente considerada a primeira axiomatização publicada da teoria dos conjuntos, Zermelo introduziu o que hoje reconhecemos como este axioma.
Nesta formulação, o Axioma do Par não aparece isoladamente, mas dentro de seu Axioma II, juntamente com a existência do vazio e dos singletos \cite{zermelo1908untersuchungen}.
Na formulação de 1930, Zermelo isolou o Axioma do Par, que passou a ser o Axioma P (\emph{Axiom der Paarung}) de sua axiomatização \cite{zermelo1930grenzzahlen}.

\begin{axiom}[Axioma do Par]\label{ax:par}\index{axioma!do par}
    Para quaisquer conjuntos $x$ e $y$, existe um conjunto $z$ que os tem como elementos.
    Em símbolos:

    \[
    \forall x \forall y \exists z (x \in z \land y \in z)
    \]
\end{axiom}

Notemos que, nesta formulação, o Axioma do Par não garante a existência de um conjunto cujos \emph{únicos} elementos sejam $x$ e $y$, mas apenas de um conjunto que contenha $x$ e $y$ como elementos.
Vários textos optam por essa formulação mais forte.
Porém, na presença do Esquema da Compreensão, essa formulação mais forte não é necessária.

\begin{lemma}[$\BST$]\label{lem:existenciaDoPar}\index{existência do par}
    Para quaisquer conjuntos $x$ e $y$, existe um único conjunto $z$ tal que, para todo $w$, $w \in z$ se e somente se $w = x$ ou $w = y$.

    Em símbolos:

    \[
    \forall x \forall y \exists! z \,\forall w (w \in z \leftrightarrow (w = x \lor w = y))
    \]
\end{lemma}

\begin{proof}
    Dado $x$ e $y$, pelo Axioma do Par, existe um conjunto $z'$ tal que $x \in z'$ e $y \in z'$.
    Agora, tomemos o conjunto $z = \{ w \in z' : w = x \lor w = y \}$ (que está bem definido, conforme a discussão da Seção~\ref{sec:axiomaDaExtensaoECompreensao}).
    
    Mostremos que $z$ satisfaz a condição desejada.
    Se $w \in z$, então $w \in z'$, então $w = x$ ou $w = y$ pela definição de $z$.
    Por outro lado, se $w=x$ ou $w=y$, então, como $x\in z'$ e $y\in z'$, segue por substituição de iguais que $w\in z'$.
    Logo, pela definição de $z$, temos $w\in z$.

    A unicidade segue do Axioma da Extensão: se $z_1$ e $z_2$ satisfazem a condição, então, para todo $w$:
    \[
    w \in z_1 \leftrightarrow (w = x \lor w = y) \leftrightarrow w \in z_2
    \]

    Assim, $\forall w\,(w\in z_1 \leftrightarrow w \in z_2)$, e então $z_1 = z_2$.
\end{proof}

Com isso, podemos definir os pares não ordenados.

\begin{definition}[$\BST$]\label{def:parNaoOrdenado}\index{par!não ordenado}
    Dado $x$ e $y$, o \emph{par não ordenado} de $x$ e $y$ é o conjunto $\{ x, y \}$ definido como o único conjunto $z$ tal que, para todo $w$, $w \in z$ se e somente se $w = x$ ou $w = y$.
    
    Formalmente, conforme discutido no capítulo 1, introduzimos na linguagem um símbolo de função binário $\{ \cdot, \cdot \}$, por meio do seguinte axioma definicional.

    \[
        \forall x \forall y \forall z (z = \{ x, y \} \leftrightarrow \forall w (w \in z \leftrightarrow (w = x \lor w = y))).
    \]

    Tal introdução é legitimada pelo Lema~\ref{lem:existenciaDoPar}, que garante a existência e unicidade do conjunto $z$ para quaisquer $x$ e $y$.
\end{definition}

Também é possível definir os \emph{singletos}, ou \emph{conjuntos unitários}.

\begin{definition}[$\BST$]\label{def:singleto}\index{singleto}\index{conjunto unitário|see{singleto}}
    Dado um conjunto $x$, o \emph{singleto}, ou \emph{conjunto unitário}, de $x$ é o conjunto $\{ x \}$ definido como o único conjunto $z$ tal que, para todo $w$, $w \in z$ se e somente se $w = x$.

    \[
        \forall x \forall z (z = \{ x \} \leftrightarrow \forall w (w \in z \leftrightarrow w = x)).
    \]
\end{definition}

Como ocorreu com os pares não ordenados, a introdução formal dessa notação deve ser justificada.
No lema seguinte, demonstramos a existência de tal conjunto e deixamos a verificação da unicidade como exercício ao leitor.
\begin{lemma}[$\BST$]\index{existência de singletos}\label{lem:existenciaDeSingletos}
    Para todo conjunto $x$, existe um único conjunto $z$ tal que, para todo $w$, $w \in z$ se e somente se $w = x$.

    Em símbolos:

    \[
        \forall x \exists! z \forall w\, (w \in z \leftrightarrow w = x)
    \]
\end{lemma}
\begin{proof}
    Dado $x$, considere $z=\{x, x\}$.
    Então, para todo $w$, $w \in z$ se e somente se $w = x$ ou $w = x$, o que é equivalente a $w = x$.
    Assim, $z$ satisfaz a condição.

    A unicidade segue do Axioma da Extensão (ver Exercício~\ref{exr:unicidadeDosSingletos}).
\end{proof}

Também é possível definir o par ordenado.
Neste texto, adotaremos a definição de Kuratowski \cite{kuratowski1921notion}.
Por razões históricas, boa parte dos textos modernos de Teoria dos Conjuntos adota essa definição, embora existam outras opções igualmente válidas.
\begin{definition}[$\BST$]\label{def:parOrdenadoKuratowski}\index{par!ordenado}\index{definição!de Kuratowski}
    Dado $x$ e $y$, o \emph{par ordenado} de $x$ e $y$ é o conjunto $(x, y)$ definido como $\{ \{ x \}, \{ x, y \} \}$.

    Formalmente, introduzimos na linguagem um símbolo de função binário $(\cdot, \cdot)$, por meio do seguinte axioma definicional.
    \[
        \forall x \forall y\, \bigl((x, y) = \{ \{ x \}, \{ x, y \} \}\bigr)
    \]

    Dizemos que um conjunto $z$ é um par ordenado se existem $x$ e $y$ tais que $z = (x, y)$.
    Em símbolos:
    \[
        \forall z\, (z \text{ é um par ordenado} \leftrightarrow \exists x \exists y (z = (x, y)))
    \]
\end{definition}

A propriedade ``é um par ordenado'' é uma notação ``em texto'' para denotar um símbolo de predicado unário introduzido por definição.
Poderíamos, alternativamente, em vez de escrever ``$z$ é um par ordenado'', escrever ``$P(z)$'', onde $P$ é um símbolo de predicado unário definido por $\forall z\,P(z) \leftrightarrow \exists x \exists y (z = (x, y))$. 

A partir dessa codificação, será possível definir posteriormente produtos cartesianos, relações e funções inteiramente em termos de conjuntos.

Insistimos no fato de que a definição de par ordenado não consiste em uma afirmação ontológica sobre a natureza do que deveria ser a noção de par ordenado, mas sim em uma definição formal por meio da qual se constroem objetos que satisfazem propriedades específicas.
No caso presente, a propriedade fundamental é a da proposição seguinte.
\begin{proposition}[$\BST$]\label{prop:propriedadeDoParOrdenado}\index{par!ordenado!propriedade}
    Para quaisquer conjuntos $x$, $y$, $u$ e $v$, $(x, y) = (u, v)$ se e somente se $x = u$ e $y = v$.

    Em símbolos:

    \[
        \forall x \forall y \forall u \forall v\, ((x, y) = (u, v) \leftrightarrow (x = u \land y = v)).
    \]
\end{proposition}
\begin{proof}
    Sejam dados $x$, $y$, $u$ e $v$.
    É imediato da lógica de primeira ordem que $(x = u \land y = v) \rightarrow (x, y) = (u, v)$.
    Para a outra direção, suponha que $(x, y) = (u, v)$.
    Pela definição de par ordenado, temos

    \[
    \{ \{ x \}, \{ x, y \} \} = \{ \{ u \}, \{ u, v \} \}
    \]
    Chamemos $A=\{\{ x \}, \{ x, y \} \}$ e $B=\{ \{ u \}, \{ u, v \} \}$.

    Como $\{x, y\}\in B$, então $\{x, y\} = \{u\}$ ou $\{x, y\} = \{u, v\}$.

    \textbf{Caso 1:} $\{x, y\} = \{u\}$.

    Neste caso, $x \in \{u\}$ e $y \in \{u\}$, então $x = u$ e $y = u$, e, portanto, $x = u$.
    Além disso, $\{u, v\} \in A$, então $\{u, v\} = \{x\}$ ou $\{u, v\} = \{x, y\}$.
    Em qualquer caso, $v=u=x=y$, e, portanto, segue a tese.


    \textbf{Caso 2:} $\{x, y\} \neq \{u\}$.

    Neste caso, como $\{u\} \in A$ e $\{u\}\neq \{x, y\}$, então $\{u\} = \{x\}$, e, portanto, $u = x$.
    
    Como $\{x, y\} \in B$ e $\{x, y\} \neq \{u\}$, temos $\{x, y\} = \{u, v\}$.
    Assim, $y=u$ ou $y=v$.
    Se $y=u$, então $\{x, y\} = \{u\}$, o que é falso.
    Logo, $y=v$ e segue a tese.
\end{proof}


\section*{Exercícios}
Os exercícios a seguir devem ser resolvidos usando apenas os axiomas da extensão, da compreensão e do par.
\begin{exercise}
    Prove que, para quaisquer conjuntos $x, y, u, v$:
    \begin{enumerate}[label=(\alph*)]
        \item $\{ x, y \} = \{ y, x \}$.
        \item $\{ x, x \} = \{ x \}$.
        \item $\{ x \} = \{ y \}$ se e somente se $x = y$.
        \item $\{x, y\}=\{u\}$ se e somente se $x=y=u$.
        \item $\{ x, y \} = \{ u, v \}$ se e somente se ($x = u$ e $y = v$) ou ($x = v$ e $y = u$).
        \item Com a definição de Kuratowski, $(x, x)=\{\{x\}\}$.
    \end{enumerate}
\end{exercise}
\begin{exercise}
    Prove que, para quaisquer conjuntos $x$ e $y$,
    \[
        (x,y)=(y,x)
    \]
    se e somente se $x=y$.
\end{exercise}
\begin{exercise}\label{exr:unicidadeDosSingletos}
    Justifique mais detalhadamente a unicidade dos singletos (Lema~\ref{lem:existenciaDeSingletos}).
\end{exercise}
\begin{exercise}
    Defina, para todo $x, y$, o conjunto $P(x, y) = \{\{\{x\}, \emptyset\}, \{\{y\}\}\}$. Prove que $P(x, y) = P(u, v)$ se e somente se $x = u$ e $y = v$.

    Tal definição é uma alternativa à definição de par ordenado de Kuratowski, e é conhecida como definição de Wiener \cite{wiener1914simplification}.
\end{exercise}
\begin{exercise}
    Defina outra definição alternativa para pares ordenados que satisfaça a propriedade de que $(x, y) = (u, v)$ se e somente se $x = u$ e $y = v$.
\end{exercise}
\begin{exercise}
    Prove que, se $z$ é um par ordenado, então existem únicos conjuntos $x$ e $y$ tais que $z=(x,y)$.

    Formalmente, prove que
    \[
        \forall z\, (z \text{ é um par ordenado} \rightarrow \exists! x \exists y (z = (x, y))),
    \]
    \[
        \forall z\, (z \text{ é um par ordenado} \rightarrow \exists! y \exists x (z = (x, y))).
    \]
\end{exercise}

\begin{exercise}\label{exr:compreensaoAutorreferente}
    Mostre que, no esquema da compreensão, é essencial que a variável $y$ que fará o papel do conjunto resultante da compreensão não apareça na fórmula $\phi$ da seguinte forma.

    Considere a fórmula $\phi(x, y)$ dada por $x \notin y$.

    Mostre que o seguinte é falso:
    \[
    \forall A \exists y \forall x (x \in y \leftrightarrow (x \in A \land x \notin y)).
    \]

    Este exercício nos mostra que tentativas de construções autorreferentes, como $\{x \in y: \phi(x, y)\}$ podem ser problemáticas em Teoria dos Conjuntos.
\end{exercise}
